\subsection{Eigenvalue and Singular-Value Interlacing}

Throughout this subsection, all matrix dimensions are positive integers.

\begin{theorem} \label{thm:rank-one-psd-eigenvalue-interlacing} \hlt{Rank-One Positive Perturbation Interlacing}\\
Let $A,E\in\operatorname{Mat}(n,n,\mathbb{C})$ be Hermitian, with $E\succeq0$ and $\operatorname{rank}(E)\leq1$.
\begin{enumerate}[label=(\roman*)]
\setlength{\itemsep}{0pt}
\item $(\forall\,i\in J_n)\,\lambda_i(A)\leq\lambda_i(A+E)$.
\item $(\forall\,i\in J_{n-1})\,\lambda_i(A+E)\leq\lambda_{i+1}(A)$.
\item $(\forall\,i\in J_n)\,\lambda_i(A-E)\leq\lambda_i(A)$.
\item $(\forall\,i\in J_{n-1})\,\lambda_i(A)\leq\lambda_{i+1}(A-E)$.
\end{enumerate}
\begin{proof}
Fix ordered orthonormal eigenbases $(\mathbf{u}_j)_{j\in J_n}$ of $A$ and $(\mathbf{v}_j)_{j\in J_n}$ of $A+E$ by Theorem \ref{thm:complex-spectral-theorem} and Definition \ref{def:increasing-hermitian-eigenvalue-ordering}(i).
\begin{enumerate}[label=(\roman*)]
\setlength{\itemsep}{0pt}
\item Fix $i\in J_n$, and set $\mathcal{U}=\operatorname{span}(\mathbf{u}_i,\ldots,\mathbf{u}_n)$ and $\mathcal{V}=\operatorname{span}(\mathbf{v}_1,\ldots,\mathbf{v}_i)$.\\
Fix a unit $\mathbf{x}\in\mathcal{U}\cap\mathcal{V}$ by Corollary \ref{cor:large-subspaces-intersect}(i), normalization, and $\dim(\mathcal{U})+\dim(\mathcal{V})=n+1$.\\
The chain $\lambda_i(A)\leq\mathbf{x}^*A\mathbf{x}\leq\mathbf{x}^*(A+E)\mathbf{x}\leq\lambda_i(A+E)$ holds by Theorem \ref{thm:rayleigh-quotient-spectral-interval} and $E\succeq0$.
\item Fix $i\in J_{n-1}$, and set $\mathcal{U}=\operatorname{span}(\mathbf{u}_1,\ldots,\mathbf{u}_{i+1})$, $\mathcal{V}=\operatorname{span}(\mathbf{v}_i,\ldots,\mathbf{v}_n)$, and $\mathcal{W}=\ker(E)$.\\
The bounds $\dim(\mathcal{U}\cap\mathcal{V})\geq2$ and $\dim(\mathcal{W})\geq n-1$ hold by Corollaries \ref{cor:large-subspaces-intersect}(ii) and \ref{cor:matrix-rank-nullity}(i), respectively.\\
Fix a unit $\mathbf{x}\in\mathcal{U}\cap\mathcal{V}\cap\mathcal{W}$ by Corollary \ref{cor:large-subspaces-intersect}(i) and normalization.\\
The chain $\lambda_i(A+E)\leq\mathbf{x}^*(A+E)\mathbf{x}=\mathbf{x}^*A\mathbf{x}\leq\lambda_{i+1}(A)$ holds by Theorem \ref{thm:rayleigh-quotient-spectral-interval}.
\item $(\forall\,i\in J_n)\,\lambda_i(A-E)\leq\lambda_i(A)$ by part (i) with $A-E$ in place of $A$.
\item $(\forall\,i\in J_{n-1})\,\lambda_i(A)\leq\lambda_{i+1}(A-E)$ by part (ii) with $A-E$ in place of $A$.
\end{enumerate}
\hlt{Structural Commentary}: Two flag intersections prove the $A+E$ bounds; substitution transfers them to $A-E$.\\
\hlt{Key Insight}: Membership in $\ker(E)$ removes the perturbation from the upper comparison.
\end{proof}
\end{theorem}

\begin{corollary} \label{cor:rank-two-one-positive-one-negative-interlacing} \hlt{Signed Rank-Two Perturbation Interlacing}\\
Let $A,H\in\operatorname{Mat}(n,n,\mathbb{C})$ be Hermitian, with $n\geq2$ and $\operatorname{rank}(H)\leq2$. Suppose that $H$ has at most one positive and one negative eigenvalue.
\begin{enumerate}[label=(\roman*)]
\setlength{\itemsep}{0pt}
\item $(\forall\,i\in J_n\setminus\{1\})\,\lambda_{i-1}(A)\leq\lambda_i(A+H)$.
\item $(\forall\,i\in J_{n-1})\,\lambda_i(A+H)\leq\lambda_{i+1}(A)$.
\end{enumerate}
\begin{proof}
Fix $H=Q\operatorname{diag}(\mu_1,\ldots,\mu_n)Q^*$ with $Q$ unitary and $\mu_1,\ldots,\mu_n\in\mathbb{R}$ by Theorem \ref{thm:complex-spectral-theorem}.\\
Set $\mu_j^+\defeq\max\{\mu_j,0\}$ and $\mu_j^-\defeq\max\{-\mu_j,0\}$ for $j\in J_n$, and set $E_\pm\defeq Q\operatorname{diag}(\mu_1^\pm,\ldots,\mu_n^\pm)Q^*$.\\
The relations $H=E_+-E_-$, $E_\pm\succeq0$, and $\operatorname{rank}(E_\pm)\leq1$ hold by construction and the sign hypothesis.\\
Set $B\defeq A-E_-$, so $A=B+E_-$ and $A+H=B+E_+$.
\begin{enumerate}[label=(\roman*)]
\setlength{\itemsep}{0pt}
\item $(\forall\,i\in J_n\setminus\{1\})\,\lambda_{i-1}(A)\leq\lambda_i(B)\leq\lambda_i(A+H)$ by Theorem \ref{thm:rank-one-psd-eigenvalue-interlacing}(ii),(i).
\item $(\forall\,i\in J_{n-1})\,\lambda_i(A+H)\leq\lambda_{i+1}(B)\leq\lambda_{i+1}(A)$ by Theorem \ref{thm:rank-one-psd-eigenvalue-interlacing}(ii),(i).
\end{enumerate}
\hlt{Structural Commentary}: The signed spectral split makes $A$ and $A+H$ positive rank-one updates of a common base.\\
\hlt{Key Insight}: Comparison through $B$ aligns the lower and upper index shifts.
\end{proof}
\end{corollary}

\begin{theorem} \label{thm:cauchy-principal-submatrix-interlacing} \hlt{Cauchy Interlacing: One Principal Deletion}\\
Let $A\in\operatorname{Mat}(n+1,n+1,\mathbb{C})$ be Hermitian, and let $B\in\operatorname{Mat}(n,n,\mathbb{C})$ be its leading principal submatrix.
\begin{enumerate}[label=(\roman*)]
\setlength{\itemsep}{0pt}
\item $(\forall\,i\in J_n)\,\lambda_i(A)\leq\lambda_i(B)$.
\item $(\forall\,i\in J_n)\,\lambda_i(B)\leq\lambda_{i+1}(A)$.
\end{enumerate}
\begin{proof}
Fix ordered orthonormal eigenbases $(\mathbf{u}_j)_{j\in J_{n+1}}$ of $A$ and $(\mathbf{v}_j)_{j\in J_n}$ of $B$ by Theorem \ref{thm:complex-spectral-theorem} and Definition \ref{def:increasing-hermitian-eigenvalue-ordering}(i).\\
Define $\iota:\mathbb{C}^n\to\mathbb{C}^{n+1}$ by $\iota(\mathbf{z})=(\mathbf{z},0)$. The identities $\norm{\iota(\mathbf{z})}_2=\norm{\mathbf{z}}_2$ and $\iota(\mathbf{z})^*A\iota(\mathbf{z})=\mathbf{z}^*B\mathbf{z}$ hold by coordinate multiplication.
\begin{enumerate}[label=(\roman*)]
\setlength{\itemsep}{0pt}
\item Fix $i\in J_n$, and set $\mathcal{S}=\operatorname{span}(\iota(\mathbf{v}_1),\ldots,\iota(\mathbf{v}_i))$ and $\mathcal{T}=\operatorname{span}(\mathbf{u}_i,\ldots,\mathbf{u}_{n+1})$.\\
Fix a unit $\mathbf{x}=\iota(\mathbf{z})\in\mathcal{S}\cap\mathcal{T}$ by Corollary \ref{cor:large-subspaces-intersect}(i), normalization, and $\dim(\mathcal{S})+\dim(\mathcal{T})=n+2$.\\
The chain $\lambda_i(A)\leq\mathbf{x}^*A\mathbf{x}=\mathbf{z}^*B\mathbf{z}\leq\lambda_i(B)$ holds by Theorem \ref{thm:rayleigh-quotient-spectral-interval}.
\item Fix $i\in J_n$, and set $\mathcal{S}=\operatorname{span}(\iota(\mathbf{v}_i),\ldots,\iota(\mathbf{v}_n))$ and $\mathcal{T}=\operatorname{span}(\mathbf{u}_1,\ldots,\mathbf{u}_{i+1})$.\\
Fix a unit $\mathbf{x}=\iota(\mathbf{z})\in\mathcal{S}\cap\mathcal{T}$ by Corollary \ref{cor:large-subspaces-intersect}(i), normalization, and $\dim(\mathcal{S})+\dim(\mathcal{T})=n+2$.\\
The chain $\lambda_i(B)\leq\mathbf{z}^*B\mathbf{z}=\mathbf{x}^*A\mathbf{x}\leq\lambda_{i+1}(A)$ holds by Theorem \ref{thm:rayleigh-quotient-spectral-interval}.
\end{enumerate}
\hlt{Structural Commentary}: The embedding preserves quadratic forms; complementary spectral flags establish both inequalities.\\
\hlt{Key Insight}: One coordinate deletion leaves the dimension sum $n+2$.
\end{proof}
\end{theorem}

\begin{corollary} \label{cor:cauchy-interlacing-general-principal-submatrix} \hlt{Cauchy Interlacing for Principal Submatrices}\\
Let $m\in\mathbb{Z}_{>0}$, and let $B\in\operatorname{Mat}(n,n,\mathbb{C})$ be a principal submatrix of a Hermitian $A\in\operatorname{Mat}(n+m,n+m,\mathbb{C})$.
\begin{enumerate}[label=(\roman*)]
\setlength{\itemsep}{0pt}
\item $(\forall\,i\in J_n)\,\lambda_i(A)\leq\lambda_i(B)$.
\item $(\forall\,i\in J_n)\,\lambda_i(B)\leq\lambda_{i+m}(A)$.
\end{enumerate}
\begin{proof}
Choose a permutation matrix $P$ such that $B$ is the leading principal submatrix of $\widetilde{A}\defeq P^*AP$.\\
The matrices $A$ and $\widetilde{A}$ have the same ordered eigenvalues by Theorem \ref{thm:similarity-preserves-characteristic-polynomial}(i).\\
For $r\in\{0,\ldots,m\}$, let $A^{(r)}$ be the leading principal submatrix of $\widetilde{A}$ of size $n+m-r$; thus $A^{(0)}=\widetilde{A}$ and $A^{(m)}=B$.
\begin{enumerate}[label=(\roman*)]
\setlength{\itemsep}{0pt}
\item $(\forall\,i\in J_n)\,\lambda_i(A^{(0)})\leq\lambda_i(A^{(1)})\leq\cdots\leq\lambda_i(A^{(m)})$ by successive applications of Theorem \ref{thm:cauchy-principal-submatrix-interlacing}(i).
\item $(\forall\,i\in J_n)\,\lambda_i(A^{(m)})\leq\lambda_{i+1}(A^{(m-1)})\leq\cdots\leq\lambda_{i+m}(A^{(0)})$ by successive applications of Theorem \ref{thm:cauchy-principal-submatrix-interlacing}(ii).
\end{enumerate}
\hlt{Structural Commentary}: Permutation makes $B$ leading; chained one-coordinate deletions accumulate the upper-index shift.\\
\hlt{Key Insight}: Each deletion raises the upper comparison index by one.
\end{proof}
\end{corollary}

\begin{corollary} \label{cor:singular-values-interlace-after-deleting-row-or-column} \hlt{Singular Values Interlace After Deleting a Row or Column}\\
Let $n\geq2$, and let $B$ be obtained from $A\in\operatorname{Mat}(n,n,\mathbb{C})$ by deleting one row or one column.\\
With singular values ordered decreasingly, $(\forall\,i\in J_{n-1})\,\sigma_i(A)\geq\sigma_i(B)\geq\sigma_{i+1}(A)$.
\begin{proof}
\begin{enumerate}[label=(\roman*)]
\setlength{\itemsep}{0pt}
\item \hlt{Row Gram Reduction}: In the row case, choose a permutation matrix $P$ with $B$ as the first $n-1$ rows of $PA$. Set $G\defeq(PA)(PA)^*$.\\
The matrix $BB^*$ is its leading principal submatrix, and $PA$ and $A$ have the same singular values by Proposition \ref{prop:unitary-invariance-of-singular-values}(i).
\item \hlt{Index Identification}: $(\forall\,j\in J_n)\,\lambda_j(G)=\sigma_{n-j+1}(A)^2$ by Proposition \ref{prop:basic-singular-value-identities}(i) and Definition \ref{def:matrix-singular-values}.\\
The identity $(\forall\,j\in J_{n-1})\,\lambda_j(BB^*)=\sigma_{n-j}(B)^2$ holds by Proposition \ref{prop:basic-singular-value-identities}(i) and Definition \ref{def:matrix-singular-values}.
\item \hlt{Cauchy Transfer}: Fix $i\in J_{n-1}$ and set $j\defeq n-i$. The chain $\sigma_{i+1}(A)^2\leq\sigma_i(B)^2\leq\sigma_i(A)^2$ holds by Theorem \ref{thm:cauchy-principal-submatrix-interlacing}(i)--(ii).\\
The row-deletion claim follows by taking nonnegative square roots.
\item \hlt{Column Reduction}: If $B$ is obtained by deleting a column, then $B^*$ is obtained from $A^*$ by deleting a row.\\
The stated inequalities hold by the row-deletion claim and Proposition \ref{prop:basic-singular-value-identities}(iii).
\end{enumerate}
\hlt{Structural Commentary}: A principal Gram block proves row deletion; adjoints transfer the result to columns.\\
\hlt{Key Insight}: $j=n-i$ matches increasing Gram eigenvalues with decreasing singular values.
\end{proof}
\end{corollary}

Figure \ref{fig:eigenvalue-interlacing-ladder} compares the interlacing windows from Theorem \ref{thm:rank-one-psd-eigenvalue-interlacing} and Corollary \ref{cor:cauchy-interlacing-general-principal-submatrix}.

\begin{illustrationfigure}[H]
  \centering
  \begin{tikzpicture}[
  >=Stealth,
  scale=1,
  axis/.style={->, gray!65, line width=0.65pt},
  adot/.style={circle, fill=blue!65!black, inner sep=1.8pt},
  bdot/.style={circle, fill=red!70!black, inner sep=1.8pt},
  cdot/.style={circle, fill=green!45!black, inner sep=1.8pt},
  guide/.style={dashed, gray!55, line width=0.55pt},
  lab/.style={font=\small, align=center, fill=white, inner sep=1pt},
  tiny/.style={font=\small, align=center, fill=white, inner sep=1pt}
]
  \draw[axis] (0,1.05) -- (3.55,1.05) node[right, lab] {$A+E$};
  \draw[axis] (0,0) -- (3.55,0) node[right, lab] {$A$};

  \foreach \x in {0.85,2.55} {
    \node[adot] at (\x,0) {};
    \draw[guide] (\x,1.22) -- (\x,-0.24);
  }
  \node[bdot] at (1.45,1.05) {};
  \draw[<->, blue!55!black] (0.85,0.34) -- (2.55,0.34);
  \node[tiny, below=5pt] at (0.85,0) {$\lambda_i(A)$};
  \node[tiny, below=5pt] at (2.55,0) {$\lambda_{i+1}(A)$};
  \node[tiny, above=5pt] at (1.45,1.05) {$\lambda_i(A+E)$};

  \draw[axis] (5.1,1.05) -- (8.95,1.05) node[right, lab] {$B$};
  \draw[axis] (5.1,0) -- (8.95,0) node[right, lab] {$A$};

  \foreach \x in {6.0,7.9} {
    \node[adot] at (\x,0) {};
    \draw[guide] (\x,1.22) -- (\x,-0.24);
  }
  \node[cdot] at (6.65,1.05) {};
  \draw[<->, green!45!black] (6.0,0.34) -- (7.9,0.34);
  \node[tiny, below=5pt] at (6.0,0) {$\lambda_i(A)$};
  \node[tiny, below=5pt] at (7.9,0) {$\lambda_{i+m}(A)$};
  \node[tiny, above=5pt] at (6.65,1.05) {$\lambda_i(B)$};
\end{tikzpicture}

  \caption{$\lambda_i(A+E)$ lies in $[\lambda_i(A),\lambda_{i+1}(A)]$, while $\lambda_i(B)$ lies in $[\lambda_i(A),\lambda_{i+m}(A)]$.}
  \label{fig:eigenvalue-interlacing-ladder}
\end{illustrationfigure}

\begin{theorem} \label{thm:schur-majorization-eigenvalues-by-diagonal} \hlt{Schur Majorization of the Diagonal}\\
Let $A\in\operatorname{Mat}(n,n,\mathbb{C})$ be Hermitian, and let $d_1\leq\cdots\leq d_n$ be its diagonal entries arranged increasingly.
\begin{enumerate}[label=(\roman*)]
\setlength{\itemsep}{0pt}
\item $(\forall\,k\in J_{n-1})\,\sum\limits_{i=1}^{k}\lambda_i(A)\leq\sum\limits_{i=1}^{k}d_i$.
\item $\sum\limits_{i=1}^{n}\lambda_i(A)=\sum\limits_{i=1}^{n}d_i$.
\end{enumerate}
\begin{proof}
Choose a permutation matrix $P$ such that the diagonal of $\widetilde{A}\defeq P^*AP$ is $(d_1,\ldots,d_n)$.\\
The matrices $A$ and $\widetilde{A}$ have the same ordered eigenvalues by Theorem \ref{thm:similarity-preserves-characteristic-polynomial}(i).\\
Set $A_k$ to be the leading $k\times k$ principal submatrix of $\widetilde{A}$ for $k\in J_n$.
\begin{enumerate}[label=(\roman*)]
\setlength{\itemsep}{0pt}
\item Fix $k\in J_{n-1}$. The inequalities $(\forall\,i\in J_k)\,\lambda_i(\widetilde{A})\leq\lambda_i(A_k)$ hold by Corollary \ref{cor:cauchy-interlacing-general-principal-submatrix}(i).\\
The identity $\sum\limits_{i=1}^{k}\lambda_i(A_k)=\operatorname{tr}(A_k)=\sum\limits_{i=1}^{k}d_i$ holds by Proposition \ref{prop:trace-determinant-from-eigenvalues}(i) and Definition \ref{def:matrix-trace-and-determinant}.\\
Thus $\sum\limits_{i=1}^{k}\lambda_i(A)\leq\sum\limits_{i=1}^{k}d_i$ by the preceding inequalities and the common ordered spectrum of $A$ and $\widetilde{A}$.
\item $\sum\limits_{i=1}^{n}\lambda_i(A)=\operatorname{tr}(A)=\sum\limits_{i=1}^{n}d_i$ by Proposition \ref{prop:trace-determinant-from-eigenvalues}(i) and Definition \ref{def:matrix-trace-and-determinant}.
\end{enumerate}
\hlt{Structural Commentary}: Cauchy bounds each leading-block eigenvalue before trace sums the inequalities.\\
\hlt{Key Insight}: Ordering the diagonal makes each partial sum a principal-block trace.
\end{proof}
\end{theorem}

\begin{proofmethod} \label{meth:eigenvalue-interlacing-method} \hlt{Eigenvalue Interlacing Proof Method}\\
Given a proposed Hermitian eigenvalue interlacing inequality, construct a proof route as follows:
\begin{enumerate}[label=(\roman*)]
\setlength{\itemsep}{0pt}
\item order every Hermitian spectrum increasingly and identify the requested lower and upper indices;
\item identify the lost or unaffected directions through perturbation rank, a kernel, or deleted coordinates;
\item choose lower and upper spectral flags whose dimensions force a nonzero intersection by Corollary \ref{cor:large-subspaces-intersect};
\item normalize an intersection vector and certify the target eigenvalue chain by Theorem \ref{thm:rayleigh-quotient-spectral-interval}.
\end{enumerate}
\end{proofmethod}
